Conditional on \(\mathrm{thm:partition-uniqueness}\), let \(\pi^\star(P)\) denote the common partition of all \(R\in\mathfrak A(P)\).  Every such representation then lies in the supplied-partition fiber \(\mathfrak G_{\pi^\star(P)}(P)\).  The cited Göbler--Windisch--Drton fixed-partition theorem implies that all their labelled block graphs coincide; call the common graph \(G^\star(P)\).  Nonemptiness of \(\mathfrak A(P)\) then gives both inclusions
\[
 \mathcal T(P)\subseteq\{(\pi^\star(P),G^\star(P))\}
 \quad\text{and}\quad
 \{(\pi^\star(P),G^\star(P))\}\subseteq\mathcal T(P),
\]
the second by choosing any \(R\in\mathfrak A(P)\).  Thus the graph step is complete and the only unresolved premise of the headline singleton formula is partition uniqueness.